% Pursuing Stacks (cotes 134-2 to 134-8): archive cote 134-3, site page 90 (PDF page 91).
% notes for ch. II. Draft compared: ps2.tex:175-184 ; ps2.tex:301-310 ; ps2.tex:328-349 ; ps2.tex:506-516.
% AI-transcribed by Claude Opus 5.5 (claude-opus-5-5, adaptive thinking, effort high), Message Batches API, batch ps1 (msgbatch_01Rv5hCaNTtUt4JHx1yqoHPL), 2026-11-12; not yet checked by eye.
\scanpage{134-3/90}
II ii

that the homotopy types modelled by two complexes which correspond should
be canonically isomorphic. In terms of the theory of \hand{``}test-functors\hand{''}which
is to follow,  this amounts essentially to the question of finding an
explicit test functor \hand{$\Delta \to \widehat{\bigcirc}$} (and\hand{/}similarly in the five other cases).
This reminds me of the problem of giving handy existence and unicity
theorems for test functors,\xout{\unsure{f}}dwelled upon somewhat in section 90, but by
no means settled yet ...

($^5$)   Compare with my letter to Breen ( I App.7).

($^6$)   Compare with (I App.2).

($^7$)   The definition of the nerve functor,  realizing a full embedding
of (Cat) into \hand{$\Delta$}$\hat{}$, is given for the first time \inserted{(I believe),} in a Bourbaki talk of mine
\margin{X} (            ),  which at that time was'nt concerned at all with the topo-%
logical interpretation of small categories,  but with \xout{\unsure{that}} some rather
formal prerequisites for the operation of "passage to quotient" of an ob-%
ject by a "pre-equivalence relation",  in any category (with a view of
applying it \xout{\unsure{to the}} in the category of schemes ...).

(\over{\illegible{1}}{$^8$})   This is well known to be false,  and the mistake is corrected
later (section 24).  The next statement,  about (Cat) modelizing the homo-%
topy category,  is true however and has been fairly well known  (due to
Fritsch and Latch, Ronnie Brown tells me).  In these notes,  it is going
to come out in section 34,  under the form "\hand{$\Delta$} is a test category",
using the general criterion for test categories obtained in section 31.

($^9$)   As seen a little later,  this \inserted{tentative ,} description of a suitable class of
cofibrations is unadequate (see  beginning of section 30).  An adequate
description of a closed model structure in (Cat) was found by \uline{Thomason}
(see comments on this in section 87).

($^{10}$)   Dispite this "feeling",  it turns out that the whole first volume
of "Pursuing Stacks" is spent on pursuing this general question o\over{r}{f}
understanding homotopy models in general,  \xout{rather than} while the stacks
are being forgotten  for a thousand pages or so ...

($^{11}$)   Finally,  it turns out more suitable still to replace this cate-%
gory by the \xout{still} smaller one of standard hemispheres \hand{$\bigcirc$},  compare ($^4$).

($^{12}$)   From a letter of Benabou which I got later,  I understand that he
got discouraged from pursuing his topological approach to n-categories,
because of the complete lack of interest of anybody else in this approach.
I was glad to hear from him last year that my present work has stimulated
him into starting a seminar on this topic of stacks,  using my notes as a
leading thread.

\mspage{90}
